\documentclass[../../../include/open-logic-section]{subfiles}

\begin{document}

\olfileid{sth}{ord-arithmetic}{using-addition}
\olsection{运用序数加法}

利用序数加法，我们可以显式计算各种集合的秩（按
\olref[spine][rank]{defnsetrank} 的定义）：


\begin{lem}\ollabel{rankcomputation}
若 $\setrank{A} = \alpha$ 且 $\setrank{B} = \beta$，则：
\begin{enumerate}
	\item\ollabel{exrankpow} $\setrank{\Pow{A}} = \alpha\ordplus 1$
	\item\ollabel{exrankpair} $\setrank{\{A, B\}} = \max(\alpha,
	\beta) \ordplus 1$
	\item\ollabel{exrankcup} $\setrank{A \cup B} = \max(\alpha,
	\beta)$
	\item\ollabel{exranktuple} $\setrank{\tuple{A,B}} = \max(\alpha,
	\beta) \ordplus  2$
	\item\ollabel{exranktimes} $\setrank{A \times B} \leq \max(\alpha,
	\beta) \ordplus  2$
	\item\ollabel{exrankunion} 当 $\alpha$ 为 0 或极限序数时，
	$\setrank{\bigcup A} = \alpha$；当 $\alpha = \gamma\ordplus 1$ 时，
	$\setrank{\bigcup A} = \gamma$。
\end{enumerate}
\end{lem}

\begin{proof}
在整个证明中，我们将多次引用 \olref[spine][rank]{ranksupstrict}。

\emph{\olref{exrankpow} 的证明。} 若 $x \subseteq A$，则 $\setrank{x} \leq
\setrank{A}$。所以 $\setrank{\Pow{A}} \leq \alpha \ordplus  1$。特别地，由于 $A
\in \Pow{A}$，可得 $\setrank{\Pow{A}} = \alpha \ordplus  1$。

\emph{\olref{exrankpair} 的证明。} 由 \olref[spine][rank]{ranksupstrict} 即得。

\emph{\olref{exrankcup} 的证明。} 由 \olref[spine][rank]{ranksupstrict} 即得。

\emph{\olref{exranktuple} 的证明。} 两次应用 \olref{exrankpair}。

\emph{\olref{exranktimes} 的证明。} 注意到 $A \times B \subseteq
\Pow{\Pow{A \cup B}}$，并引用 \olref{exranktuple}。

\emph{\olref{exrankunion} 的证明。} 若 $\alpha = \gamma\ordplus 1$，则存在
$c \in A$ 使得 $\setrank{c} = \gamma$，并且 $A$ 中没有元素具有更高的秩；
因此 $\setrank{\bigcup A} = \gamma$。若 $\alpha$ 是极限序数，则 $A$ 包含的元素的秩可以任意接近（但严格小于）$\alpha$，于是 $\bigcup A$ 也包含秩可以任意接近（但严格小于）$\alpha$ 的元素，因此 $\setrank{\bigcup A} = \alpha$。
\end{proof}


\noindent
我们将此留作练习，以说明为何 \olref{exranktimes} 涉及的是一个\emph{不等}式。


\begin{prob}
构造集合 $A$ 和 $B$ 使得 $\setrank{A \times B}=
\max(\setrank{A}, \setrank{B})$。构造集合 $A$ 和 $B$ 使得
$\setrank{A \times B}\max(\setrank{A}, \setrank{B}) \ordplus  2$。是否
还存在其他可能的秩？
\end{prob}



现在我们还可以证明，将一个序数描述为“有限”或“无穷”的几种合理概念是相互一致的：


\begin{lem}\ollabel{ordinfinitycharacter}
对任意序数 $\alpha$，下列条件等价：
\begin{enumerate}
	\item\ollabel{ord:notinomega} $\alpha\notin \omega$，即 $\alpha$ 不是自然数
	\item\ollabel{ord:omegaplus} $\omega \leq \alpha$
	\item\ollabel{ord:oneplus} $1 \ordplus  \alpha = \alpha$
	%\alpha \approx \alpha \ordplus 1$
	\item\ollabel{ord:plusone} $\alpha \approx \alpha\ordplus 1$，
	即 $\alpha$ 与 $\alpha\ordplus 1$ 等势
	\item\ollabel{ord:infinite} $\alpha$ 是 Dedekind（戴德金）无穷的
	\end{enumerate}
\end{lem}


\noindent
因此，我们有了五种可证明等价的方式来理解一个序数何时为（非）有限的。


\begin{proof}
\emph{\olref{ord:notinomega} 推出 \olref{ord:omegaplus}。} 由
三分律可得。

\emph{\olref{ord:omegaplus} 推出 \olref{ord:oneplus}。} 取定
$\alpha \geq \omega$。根据超限归纳法，存在某个最小序数 $\gamma$（可能是 $0$），使得存在一个极限序数
$\beta$ 满足 $\alpha = \beta \ordplus \gamma$。现在：
\[
	1 \ordplus \alpha =  
	1 \ordplus (\beta \ordplus \gamma) = 
	(1 \ordplus \beta) \ordplus \gamma =  
	\supstrict_{\delta < \beta} (1 \ordplus  \delta) \ordplus  \gamma = 
	\beta \ordplus  \gamma = 
	\alpha.
\]
\emph{\olref{ord:oneplus} 推出 \olref{ord:plusone}。} 显然存在一个双射 $f \colon (\alpha \disjointsum 1) \to (1
\disjointsum \alpha)$。若 $1 \ordplus \alpha = \alpha$，则存在一个同构 $g \colon (1 \disjointsum \alpha) \to \alpha$。现在
考虑 $\comp{f}{g}$。

\emph{\olref{ord:plusone} 推出 \olref{ord:infinite}。} 若
$\alpha \approx \alpha \ordplus 1$，则存在一个双射 $f \colon
(\alpha \disjointsum 1) \to \alpha$。对每个 $\gamma < \alpha$ 定义 $g(\gamma) = f(\gamma, 0)$；该单射见证了 $\alpha$ 是 Dedekind 无穷的，因为 $f(0,1) \in \alpha \setminus \ran{g}$。

\emph{\olref{ord:infinite} 推出 \olref{ord:notinomega}。} 这是
\olref[z][infinity-again]{naturalnumbersarentinfinite}。
\end{proof}



\end{document}